\section{Introduction}
These are Alonso's notes for an introductory course on one step stochastic processes in continuous Time. These type of processes are often used in science to model uncertainty and variability in data collected from a given experiment or a system through observation. This field is huge and here we focus only on one type of stochastic process. These notes are intended to present the pieces of theory we will need to motivate ---and formalize with some mathematical rigor--- one-step (also known as birth-death) stochastic processes in continuous time. We assume the reader is familiar with basic calculus and has been exposed to an introductory course on probability and statistics. As a refresher, we begin with basic concepts and definitions of probability theory. Then, as a first simple example of a one-step stochastic process, we derive the Poisson distribution that results from a counting process. We derive the so-called master equation in this simple case, and show how to solve it. It follows then a general derivation of the master equation for one-step processes in one dimension, and how to find the probability distribution characterizing the stationary state. We also present basic theory on quasi-stationary distributions and first-passage times. Then, through simple examples, the essence of the Gillespie algorithm is explained as a way to build exact stochastic simulations. The final sections present examples and case studies of increasing complexity of the application of stochastic processes in continuous time in the context of population and community ecology and the dynamics of infectious diseases. Particular attention is paid, when possible, to show the connection between stochastic models and empirical data through the likelihood function.
\smallskip

We hope you will find these notes useful. 
